% !TEX program = xelatex
\documentclass[11pt,a4paper]{article}

\usepackage[margin=1.7cm]{geometry}
\usepackage{fontspec}
\usepackage{polyglossia}
\setdefaultlanguage{spanish}
\usepackage{booktabs}
\usepackage{tabularx}
\usepackage{array}
\usepackage{enumitem}
\usepackage{hyperref}
\usepackage{xcolor}
\usepackage{graphicx}
\usepackage{tikz}
\usetikzlibrary{arrows.meta, positioning, fit, calc, backgrounds}
\usepackage{fancyhdr}
\usepackage{titlesec}
\usepackage{microtype}
\usepackage{xspace}

\setmainfont{Arial}
\setsansfont{Arial}

\hypersetup{colorlinks=true, linkcolor=black, urlcolor=blue!50!black}
\pagestyle{fancy}
\fancyhf{}
\fancyfoot[C]{\thepage}
\renewcommand{\headrulewidth}{0pt}

\titleformat{\section}{\bfseries\large}{\thesection.}{0.4em}{}
\titleformat{\subsection}{\bfseries\normalsize}{\thesubsection.}{0.35em}{}
\titlespacing*{\section}{0pt}{0.7em}{0.25em}
\titlespacing*{\subsection}{0pt}{0.45em}{0.15em}
\setlist[itemize]{leftmargin=1.15em, itemsep=0.05em, topsep=0.1em}
\setlist[enumerate]{leftmargin=1.25em, itemsep=0.05em, topsep=0.1em}

\newcommand{\piroute}{$\pi(x)$\xspace}

\begin{document}

\begin{center}
{\LARGE\bfseries Pipeline experimental completo y reproducible}\\[0.25em]
{\large Proyecto de Investigación II}\\[0.45em]
\begin{tabular}{@{}p{0.22\textwidth}p{0.72\textwidth}@{}}
\bfseries Título del plan: &
Asistente de consultas académicas-administrativas con enrutamiento predictivo \\[0.15em]
\bfseries Caso: &
Facultad de Ingeniería y Negocios, UMA (San Juan de Lurigancho) \\[0.15em]
\bfseries Tipo según guía: &
LLM + RAG, herramientas sobre malla, y enrutamiento predictivo de estrategias \\
\end{tabular}
\end{center}

\vspace{0.3em}
\noindent
La clave es que otra persona pueda ejecutar el proceso y comprender cómo se pasa de los documentos originales a los resultados. Este es el diseño del flujo experimental para la tesis.

\section{Vista del pipeline (cómo se lee el documento)}

Hay dos diagramas. El primero muestra el \textbf{pipeline experimental} (de los documentos al informe). El segundo muestra la \textbf{arquitectura en tiempo de consulta} (qué hace el sistema una vez construido). Las fases~1--12 detallan cada caja.

\subsection*{Figura 1. Pipeline experimental (doce fases)}

\begin{center}
\begin{tikzpicture}[
  font=\scriptsize,
  >=Latex,
  stage/.style={draw, rounded corners=2pt, align=center, inner sep=2.8pt,
                minimum width=2.05cm, minimum height=0.95cm},
  arr/.style={->, thick},
]
% Fila 1: datos
\node[stage, fill=blue!12] (f1) {1. Definición\\experimento};
\node[stage, fill=blue!12, right=0.22cm of f1] (f2) {2. Adquisición\\de datos};
\node[stage, fill=blue!12, right=0.22cm of f2] (f3) {3. Extracción\\y limpieza};
\node[stage, fill=blue!12, right=0.22cm of f3] (f4) {4. EDA\\del corpus};
\node[stage, fill=blue!12, right=0.22cm of f4] (f5) {5. Banco de\\consultas};
\node[stage, fill=blue!12, right=0.22cm of f5] (f6) {6. Validación\\dataset};

\foreach \a/\b in {f1/f2,f2/f3,f3/f4,f4/f5,f5/f6}
  \draw[arr] (\a) -- (\b);

% Fila 2: experimento
\node[stage, fill=orange!15, below=0.55cm of f1] (f7) {7. Partición\\conjuntos};
\node[stage, fill=orange!15, right=0.22cm of f7] (f8) {8. Estrategias\\A0--A5};
\node[stage, fill=orange!15, right=0.22cm of f8] (f9) {9. Ejecución\\controlada};
\node[stage, fill=green!12, right=0.22cm of f9] (f10) {10. Evaluación\\métricas};
\node[stage, fill=green!12, right=0.22cm of f10] (f11) {11. Análisis\\estadístico};
\node[stage, fill=green!12, right=0.22cm of f11] (f12) {12. Repro.\\e informe};

\foreach \a/\b in {f7/f8,f8/f9,f9/f10,f10/f11,f11/f12}
  \draw[arr] (\a) -- (\b);

\draw[arr] (f6.south) -- ++(0,-0.22) -| (f7.north);

% Leyenda de bandas
\begin{scope}[on background layer]
  \node[fit=(f1)(f6), inner sep=4pt, fill=blue!4, rounded corners=3pt] {};
  \node[fit=(f7)(f9), inner sep=4pt, fill=orange!6, rounded corners=3pt] {};
  \node[fit=(f10)(f12), inner sep=4pt, fill=green!5, rounded corners=3pt] {};
\end{scope}

\node[font=\tiny, anchor=west, text=blue!50!black] at ($(f1.west)+(-0.05,0.72)$)
  {Datos y banco};
\node[font=\tiny, anchor=west, text=orange!70!black] at ($(f7.west)+(-0.05,-0.72)$)
  {Implementación y corridas};
\node[font=\tiny, anchor=east, text=green!50!black] at ($(f12.east)+(0.05,-0.72)$)
  {Resultados};
\end{tikzpicture}
\end{center}

{\footnotesize Figura~1. De la documentación original al informe final. Azul: preparar datos y banco. Naranja: implementar y ejecutar. Verde: medir y reportar.}

\subsection*{Figura 2. Arquitectura en tiempo de consulta}

\begin{center}
\resizebox{\textwidth}{!}{%
\begin{tikzpicture}[
  font=\scriptsize,
  >=Latex,
  box/.style={draw, rounded corners=2pt, align=center, inner sep=3pt},
  arr/.style={->, thick},
  dasharr/.style={->, thick, densely dashed},
  strat/.style={draw, rounded corners=1pt, align=center, fill=gray!8,
                minimum width=2.5cm, minimum height=0.52cm, inner sep=2pt},
]
% Columna 1-2: entrada
\node[box, fill=yellow!15] (q) at (0,2.4) {Consulta\\estudiante};
\node[box, fill=blue!10] (an) at (2.4,2.4) {Análisis\\tipo / fuentes};
\node[box, fill=orange!20] (rou) at (4.7,2.4) {Router \piroute\\(solo A5)};

% Columna estrategias (centrada, espaciada)
\node[strat] (a0) at (7.2,3.7) {A0 base};
\node[strat] (a1) at (7.2,3.05) {A1 RAG};
\node[strat] (a2) at (7.2,2.4) {A2 FT (si viable)};
\node[strat] (a3) at (7.2,1.75) {A3 RAG + tools};
\node[strat] (a4) at (7.2,1.1) {A4 compuesta fija};

% Salida
\node[box, fill=green!12] (out) at (10.0,2.4) {Respuesta\\+ fuentes / persona};
\node[box, fill=green!18] (met) at (12.5,2.4) {Métricas\\cal./lat./costo};

\draw[arr] (q) -- (an);
\draw[arr] (an) -- (rou);
\foreach \n in {a0,a1,a2,a3,a4}
  \draw[arr] (rou.east) -- ++(0.2,0) |- (\n.west);
\foreach \n in {a0,a1,a3,a4}
  \draw[arr] (\n.east) -- ++(0.2,0) |- (out.west);
\draw[arr] (a2) -- (out);
\draw[arr] (out) -- (met);

% Bypass A0--A4 (sin router)
\draw[dasharr, gray] (an.south) -- ++(0,-0.45) -| (a4.south);

% Datos debajo, con hueco claro respecto a A4 (y=1.1)
\node[box, fill=blue!8, minimum width=2.8cm] (corp) at (2.4,-0.35)
  {Corpus docs\\norma / calendario};
\node[box, fill=blue!8, minimum width=2.8cm] (mall) at (7.2,-0.35)
  {Grafo malla\\tools / reglas};
\node[box, fill=yellow!12, minimum width=2.8cm] (gold) at (11.5,-0.35)
  {Banco gold\\validado};

\draw[dasharr] (corp) |- (a1);
\draw[dasharr] (mall) -- (a3);
\draw[dasharr] (gold) |- (met);

\node[font=\tiny, align=left, text width=12.5cm] at (6.5,-1.25)
  {A5 = con router (elige A0--A4). A0--A4 fijas: flecha gris. Discontinuas: RAG, tools, evaluación.};
\end{tikzpicture}%
}
\end{center}

{\footnotesize Figura~2. Arquitectura al atender una consulta. El experimento compara A0--A5 sobre el mismo banco.}

\section{Fase 1. Definición del experimento}

Antes de programar, documentar:

\begin{itemize}
  \item Pregunta de investigación.
  \item Hipótesis.
  \item Alcance del dominio (consultas académicas-administrativas en la facultad UMA--SJL).
  \item Estrategias que se compararán (A0--A5).
  \item Métricas principales.
  \item Criterios de éxito.
  \item Limitaciones previstas.
\end{itemize}

\paragraph{Pregunta central.}
¿El routing predictivo de estrategias de IA mejora la relación entre calidad, latencia y costo frente a utilizar una estrategia fija para todas las consultas académicas-administrativas?

No hace falta suponer que el routing será superior. El experimento debe permitir que los resultados confirmen o contradigan esa hipótesis.

\section{Fase 2. Adquisición de datos}

\paragraph{Entradas.}
Documentación pública; documentación institucional autorizada; preguntas frecuentes; consultas reales anonimizadas, si se consiguen.

\paragraph{Procedimiento.}
\begin{itemize}
  \item Registrar cada fuente.
  \item Descargar los archivos permitidos.
  \item Conservar metadatos y versiones.
  \item Registrar fecha de consulta y procedencia.
  \item Documentar restricciones de uso.
  \item Separar fuentes públicas de información autorizada no pública.
\end{itemize}

\paragraph{Salida.}
Un inventario documental reproducible.

\section{Fase 3. Extracción y limpieza}

Procesar los archivos con herramientas adecuadas para PDF, HTML y documentos de texto.

\paragraph{Tareas.}
\begin{itemize}
  \item Extracción de texto.
  \item Identificación de tablas.
  \item Normalización de caracteres.
  \item Detección de duplicados.
  \item Detección de documentos vacíos o mal procesados.
  \item Registro de errores.
  \item Identificación de documentos potencialmente desactualizados.
\end{itemize}

No eliminar indiscriminadamente tablas, encabezados o referencias legales: pueden ser esenciales para responder preguntas.

\paragraph{Salida.}
Corpus procesado con trazabilidad hacia la fuente original. La malla curricular, si se usa, se estructura aparte para herramientas o reglas (prerrequisitos), no como texto mezclado sin criterio en el índice de recuperación.

\section{Fase 4. EDA del corpus}

Analizar:

\begin{itemize}
  \item Tipos de documentos.
  \item Años y versiones.
  \item Longitud de los textos.
  \item Calidad de extracción.
  \item Cobertura temática.
  \item Vacíos de información.
  \item Duplicados y posibles contradicciones.
\end{itemize}

Generar tablas y gráficos, pero acompañarlos de conclusiones que expliquen cómo afectan al sistema.

\paragraph{Salida.}
Informe de EDA documental y decisiones de preparación del corpus.

\section{Fase 5. Construcción del dataset de consultas}

Construir un banco inicial con preguntas extraídas de las FAQ, propuestas por el equipo y generadas a partir de documentos oficiales.

Cada registro debe incluir, como mínimo:

\begin{itemize}
  \item Identificador.
  \item Pregunta.
  \item Categoría.
  \item Fuente de referencia.
  \item Respuesta de referencia, cuando se haya verificado.
  \item Estado de validación.
\end{itemize}

Se puede añadir complejidad, número de fuentes requeridas, estrategia candidata y si la pregunta es respondible con el corpus disponible.

Las etiquetas deben basarse en reglas documentadas y revisión humana, no únicamente en la opinión de un modelo de IA.

\section{Fase 6. Validación del dataset}

Los validadores revisan las preguntas y sus respuestas.

\paragraph{Procedimiento.}
\begin{itemize}
  \item Comprobar que la consulta sea pertinente.
  \item Verificar la respuesta contra la documentación.
  \item Confirmar la referencia documental.
  \item Identificar ambigüedades.
  \item Marcar preguntas sin respuesta verificable.
  \item Revisar los casos que requieren varias fuentes.
\end{itemize}

Documentar los desacuerdos entre validadores y cómo se resolvieron.

\paragraph{Salida.}
Banco de evaluación validado.

\section{Fase 7. Preparación de los conjuntos de datos}

Si se entrena un router o se hace fine-tuning, separar los datos según su función. Por ejemplo: entrenamiento, validación y prueba final.

La partición debe evitar fugas de información: preguntas casi idénticas o derivadas del mismo caso no deberían repartirse indiscriminadamente entre entrenamiento y prueba.

El conjunto de prueba debe quedar separado de las decisiones de ajuste.

Si el router utiliza un clasificador entrenado, registrar sus características de entrada, etiquetas, algoritmo, hiperparámetros y semilla aleatoria.

\section{Fase 8. Implementación de las estrategias candidatas}

Implementar las configuraciones que realmente se puedan evaluar:

\begin{center}
\begin{tabular}{@{}clp{0.58\textwidth}@{}}
\toprule
\textbf{Config.} & \textbf{Nombre} & \textbf{Descripción} \\
\midrule
A0 & Modelo base & Sin recuperación documental. \\
A1 & RAG & Modelo con recuperación documental. \\
A2 & Fine-tuning & Modelo con fine-tuning, si es viable. \\
A3 & RAG + herramientas & RAG con herramientas para información estructurada. \\
A4 & Compuesta fija & Estrategia compuesta fija, si su complejidad está justificada. \\
A5 & Routing predictivo & Selecciona entre estrategias candidatas. \\
\bottomrule
\end{tabular}
\end{center}

El fine-tuning y el multiagente no tienen que incluirse a la fuerza si los recursos o los datos no permiten evaluarlos correctamente.

Para la estrategia basada en mallas, se pueden usar herramientas o reglas explícitas para consultar prerrequisitos, siempre que esa información esté disponible y pueda estructurarse correctamente.

\section{Fase 9. Ejecución controlada}

Todas las configuraciones deben evaluarse sobre el mismo conjunto de consultas.

Registrar:

\begin{itemize}
  \item Identificador de la consulta.
  \item Configuración ejecutada.
  \item Respuesta producida.
  \item Fuentes recuperadas.
  \item Latencia.
  \item Tokens consumidos, cuando corresponda.
  \item Costo estimado.
  \item Ruta seleccionada.
  \item Errores y excepciones.
\end{itemize}

Mantener constantes, en la medida de lo posible, los modelos, versiones, parámetros y condiciones de ejecución.

Si se usan servicios externos, registrar sus versiones o identificadores, parámetros y fechas de ejecución.

\section{Fase 10. Evaluación}

Evaluar los resultados desde varias dimensiones.

\paragraph{Calidad.}
Exactitud de la respuesta; corrección de la información; fidelidad a las fuentes; corrección de las citas; tasa de respuestas sin fundamento.

\paragraph{Eficiencia.}
Latencia mediana; percentil~95 de latencia; costo por consulta; tokens consumidos.

\paragraph{Routing.}
Selección de la estrategia adecuada según los criterios establecidos; errores de selección; frecuencia de uso de cada estrategia; casos en los que el sistema selecciona una estrategia innecesariamente costosa.

Definir las métricas y sus reglas de cálculo antes de ejecutar la evaluación final.

\section{Fase 11. Análisis estadístico}

Comparar las configuraciones:

\begin{itemize}
  \item En el conjunto completo.
  \item Por categoría.
  \item Por complejidad.
  \item En consultas que necesitan varias fuentes.
  \item En consultas que necesitan herramientas.
\end{itemize}

Utilizar comparaciones pareadas porque las mismas consultas se ejecutan con varias estrategias.

Cuando corresponda, reportar intervalos de confianza y pruebas estadísticas adecuadas al tipo de métrica y diseño experimental.

No limitarse a demostrar que una configuración es más rápida: analizar si mantiene una calidad aceptable.

\section{Fase 12. Reproducibilidad y documentación}

Registrar:

\begin{itemize}
  \item Versiones de dependencias.
  \item Semillas aleatorias.
  \item Configuración de modelos.
  \item Versiones del corpus.
  \item Parámetros de recuperación.
  \item Configuración del router.
  \item Fecha de ejecución.
  \item Datos de entrada utilizados.
  \item Resultados generados.
\end{itemize}

El repositorio debe permitir ejecutar el pipeline sin depender de rutas locales de una sola computadora ni de archivos que no estén documentados.

\paragraph{Flujo resumido.}
Documentación $\to$ extracción $\to$ EDA $\to$ banco de consultas $\to$ validación $\to$ preparación de datos $\to$ implementación $\to$ ejecución experimental $\to$ evaluación $\to$ análisis estadístico $\to$ informe final.

\end{document}
